% SEGMENT OLP-0656-S01
% Part: proof-theory
% Chapter: cut-elimination
% Section: ce-topmost

\documentclass[../../../include/open-logic-section]{subfiles}

\begin{document}

\olfileid{pt}{cut}{top}

\olsection{உச்சியிலுள்ள வெட்டுகளை நீக்குதல்}

\begin{defn}
\CutCS{} விதியில் முடியும் !!a{proof}~$\pi$ ஐக் கருதுவோம்:
\[
\AxiomC{}
\RightLabel{$\pi_1$}
\Deduce$\Gamma \fCenter \Delta, !A$
\AxiomC{}
\RightLabel{$\pi_2$}
\Deduce$!A, \Gamma \fCenter \Delta$
\RightLabel{\CutCS}
\BinaryInf$\Gamma \fCenter \Delta$
\DisplayProof
\]


இதில் வேறு எந்த \CutCS{} அனுமானமும் இல்லை என்க.
~$\pi$ இன் \emph{வெட்டு உயரம்}
$\cheight{\pi} = \pheight{\pi_1} + \pheight{\pi_2}$.
~$\pi$ இன் \emph{வெட்டுத் தரம்}
$\cutrank{\pi} = \depth{!A}$.
\end{defn}

% SEGMENT OLP-0656-S02
\begin{lem}\ollabel{lem:cut-adm-G3c}
\CutCS{} விதி~\Log{G3c} இல் ஏற்கத்தக்கது.
வேறு சொற்களில், ஒரே ஒரு~\CutCS{} இல் முடிந்து,
மற்ற இடங்களில்~$\Log{G3c}$ இன் விதிகளை
மட்டும் பயன்படுத்தும் !!a{proof}~$\pi$ இருந்தால்,
அதே இறுதித் தொடரணிக்கான !!a{proof}~$\pi'$
ஆனது~\Log{G3c} இல் உண்டு.
\end{lem}

\begin{proof}
வெட்டு !!{height}, வெட்டுத் தரம் ஆகிய
இரண்டின் மீது தொகுத்தறிந்து நிறுவுவோம்.
!!{proof}~$\pi$ இவ்வாறு முடிகிறது:
\[
\AxiomC{}
\RightLabel{$\pi_1$}
\Deduce$\Gamma \fCenter \Delta, !A$
\AxiomC{}
\RightLabel{$\pi_2$}
\Deduce$!A, \Gamma \fCenter \Delta$
\RightLabel{\CutCS}
\BinaryInf$\Gamma \fCenter \Delta$
\DisplayProof
\]

தொகுத்தறிதலின் அடிப்படை நேர்வு:
$\cheight{\pi} = 0$ மற்றும்
$\cutrank{\pi} = 0$.
$\cheight{\pi} = \pheight{\pi_1} + \pheight{\pi_2} = 0$
ஆகையால், $\Gamma \Sequent \Delta, !A$,
$!A, \Gamma \Sequent \Delta$
இரண்டும் அடிக்கோள்கள்.
இரண்டு வாய்ப்புகள் உள்ளன:
(a)~அவற்றில் ஒன்றிலாவது $!A$
முதன்மை வாய்பாடல்ல; அல்லது
(b)~இரண்டிலும் $!A$ முதன்மை வாய்பாடு.
இரண்டு நேர்வுகளிலும்
$\Gamma \Sequent \Delta$ ஓர்
அடிக்கோளாகவே இருக்க வேண்டும்
என்பதைக் கீழே காட்டுவோம்;
அதற்கு தொகுத்தறிதல் கருதுகோள்
தேவையில்லை.

இனி~$\pi_1$, $\pi_2$ ஆகியவை
அடிக்கோள்களா அல்லது அனுமானத்தில்
முடிகின்றனவா, அந்த அனுமானத்தில்
~$!A$ முதன்மை வாய்பாடா இல்லையா
என்பவற்றைப் பொறுத்து நேர்வுகளைப்
பிரிப்போம். A, B நேர்வுகள்
தொகுத்தறிதலின் அடிப்படை நேர்வை
உள்ளடக்கும். தொகுத்தறிதல் படியில்
$\cheight{\pi} > 0$ அல்லது
$\cutrank{\pi} > 0$
(அல்லது இரண்டும்) எனக் கொள்வோம்.
ஒரே ஒரு~\CutCS{} இல் முடியும்
!!{proof} ஒன்றில், வெட்டுத் தரம்
$= \cutrank{\pi}$ ஆகவும்
வெட்டு !!{height} ஆனது
$< \cheight{\pi}$ ஆகவும் இருந்தாலோ,
அல்லது வெட்டுத் தரம்
$< \cutrank{\pi}$ ஆக இருந்தாலோ,
அதன் இறுதி~\CutCS{} ஐ நீக்கலாம்
என்ற தொகுத்தறிதல் கருதுகோளைப்
பயன்படுத்தலாம். $\cheight{\pi} = 0$
என்ற நேர்வை, அதாவது இரண்டு
முன்கோள்களும் அடிக்கோள்களான
நேர்வை, A மற்றும் B நேர்வுகள்
உள்ளடக்குகின்றன.
$\cheight{\pi} > 0$
என்ற நேர்வுகளை C--D நேர்வுகள்
உள்ளடக்குகின்றன.

% SEGMENT OLP-0656-S03
\paragraph{நேர்வு A:}
ஏதோ ஒரு $!B$ ஆனது~$\Gamma$,
~$\Delta$ இரண்டிலும் வருகிறது.
அப்போது $!B$ அணுவாய்பாடாக
இருந்தால் $\Gamma \Sequent \Delta$
ஓர் அடிக்கோள்; இல்லையேல்
\olref[seq][ex]{prop:G3c-axioms} மூலம்
\CutCS{} இன்றி~\Log{G3c} இல்
அதற்கு !!a{proof}~$\pi'$
உண்டு. இது அடிப்படை
நேர்வின் (a) பகுதியை
உள்ளடக்குகிறது: அடிக்கோள்களான
$\Gamma \Sequent \Delta, !A$
அல்லது $!A, \Gamma \Sequent \Delta$
ஆகியவற்றில் ஒன்றிலாவது
$!A$ முதன்மை வாய்பாடல்ல எனில்,
வேறோர் அணு~$!B$
~$\Gamma$,~$\Delta$
இரண்டிலும் வரவேண்டும்.
ஒரு முன்கோள் அடிக்கோளாகவும்
அதில் $!A$ முதன்மை வாய்பாடல்லாமலும்
இருந்தால், மற்ற முன்கோள் ஒரு
விதியின் முடிவாக இருந்தாலும்,
தொகுத்தறிதல் படியையும்
இந்நேர்வு உள்ளடக்குகிறது.

\paragraph{நேர்வு B:}
இரு முன்கோள்களும்
அடிக்கோள்கள்; $!A$
இரண்டிலும் முதன்மை
வாய்பாடு, ஆகவே அணுவாய்பாடு.
இடப்புற முன்கோளான
$\Gamma \Sequent \Delta, !A$
ஐக் கருதுவோம்.
$!A$ முதன்மை வாய்பாடு
என்பதால் அது~$\Gamma$ இல்
வரவேண்டும்; இல்லையேல்
$\Gamma \Sequent \Delta, !A$
அடிக்கோளாகாது.
அதேபோல் $!A$ ஆனது
~$\Delta$ விலும் வரவேண்டும்;
இல்லையேல்
$!A, \Gamma \Sequent \Delta$
அடிக்கோளாகாது.
எனவே $!A$ அணுவாய்பாடாகி
~$\Gamma$,~$\Delta$
இரண்டிலும் வருகிறது;
அதாவது
$\Gamma \Sequent \Delta$
ஓர் அடிக்கோள்.
இது தொகுத்தறிதலின்
அடிப்படை நேர்வின் (b)
பகுதியை உள்ளடக்குகிறது.

% SEGMENT OLP-0656-S04
\begin{center}
A, B நேர்வுகளுக்கான மாற்று வழிகள்
\end{center}

\paragraph{நேர்வு A:}
வெட்டின் முன்கோள்களில் ஒன்று,
$!C$ அணுவாய்பாடாக இருக்கும்
$!C, \Gamma' \Sequent \Delta', !C$
வடிவிலான அடிக்கோள்.
$!C$ வெட்டு !!{formula}~$!A$
அல்ல எனில், $!C$ ஆனது
~$\Gamma$,~$\Delta$
இரண்டிலும் வரவேண்டும்.
அப்போது
$\Gamma \Sequent \Delta$
ஓர் அடிக்கோள்; அதையே~$\pi'$
ஆகக் கொள்ளலாம்.
$!C \ident !A$ வெட்டு
!!{formula} ஆக இருந்தால்,
இடப்புற முன்கோள்
அடிக்கோளான நேர்வில்
$!A \in \Gamma$ மற்றும்
$\Gamma = \Gamma', !A$;
வலப்புற முன்கோள்
அடிக்கோளான நேர்வில்
$!A \in \Delta$ மற்றும்
$\Delta = \Delta', !A$.
முதல் நேர்வில்~$\pi_2$
என்பது
$!A, !A, \Gamma' \Sequent \Delta$
இன் !!a{proof}; ஆகவே
\LeftR{\Contraction} இன்
ஏற்கத்தக்கதன்மையால்
(\olref[seq][inv]{prop:G3c-cont-adm})
$\Gamma \Sequent \Delta$
இன் !!a{proof} கிடைக்கும்.
இரண்டாம் நேர்வில்~$\pi_1$
என்பது
$\Gamma \Sequent \Delta', !A, !A$
இன் !!a{proof}; ஆகவே
\RightR{\Contraction} இன்
ஏற்கத்தக்கதன்மையால்
$\Gamma \Sequent \Delta$
இன் !!a{proof} கிடைக்கும்.

\paragraph{நேர்வு B:}
முன்கோள்களில் ஒன்று
$\lfalse, \Gamma' \Sequent \Delta'$
வடிவிலான அடிக்கோள்.
இடப்புற முன்கோள்
அத்தகைய அடிக்கோள் எனில்,
$\lfalse \in \Gamma$;
எனவே $\Gamma \Sequent \Delta$
உம் அடிக்கோள்.
வலப்புற முன்கோள்
அத்தகைய அடிக்கோள் எனில்,
$\lfalse \in \Gamma$
(அப்போது
$\Gamma \Sequent \Delta$
மீண்டும் அடிக்கோள்)
அல்லது $!A \ident \lfalse$.
பிந்தைய நேர்வில்~$\pi_1$
என்பது
$\Gamma \Sequent \Delta, \lfalse$
இன் !!a{proof}.
~$\pi_1$ இன் ஒவ்வொரு
தொடரணியின் முடிவுப்பக்கத்திலிருந்தும்
~$\lfalse$ இன் ஒரு
நிகழ்வை நீக்கி,
$\Gamma \Sequent \Delta$
இன் !!a{proof} ஆக
மாற்றலாம்.
(பயிற்சி: ~$\pheight{\pi_1}$
மீது தொகுத்தறிந்து
இதை உறுதிசெய்க.)

% SEGMENT OLP-0656-S05
இனி A, B நேர்வுகள்
எதுவும் பொருந்தாது என்க.
மீதமுள்ள நேர்வுகளில்
$\cheight{\pi} > 0$;
அதாவது குறைந்தது
ஒரு முன்கோளாவது ஓர்
அனுமானத்தின் முடிவாகும்.

\paragraph{C, D நேர்வுகள்: வெட்டுகளின் வரிசையை மாற்றுதல்.}
C, D நேர்வுகளுக்கெல்லாம்
ஒரே பொது வடிவம் உண்டு.
ஒரு முன்கோள்~$R$
விதியால் கிடைக்கும்
வகையில், \CutCS{} இல்
முடியும் !!a{proof}
ஒன்றிலிருந்து தொடங்குகிறோம்.
அந்த~$R$ இல் வெட்டு
!!{formula}~$!A$
முதன்மையல்ல; வேறொரு
!!{formula}~$!D$
முதன்மையானது.
அந்த !!{proof} இல்
விதி~$R$ க்குப் பிறகு
\CutCS{} வருகிறது.
அதன் மாதிரி வடிவம்:
\[
\AxiomC{}
\RightLabel{$\pi_1$}
\Deduce$\Gamma \fCenter \Delta', !D, !A$
\AxiomC{}
\RightLabel{$\pi_3$}
\Deduce$!A, \Gamma, \Pi \fCenter \Delta', \Lambda$
\RightLabel{$R$}
\UnaryInf$!A, \Gamma \fCenter \Delta', !D$
\RightSubproofLabel{$\pi_2$}
\RightLabel{\CutCS}
\BinaryInf$\Gamma \fCenter \Delta', !D$
\DisplayProof
\]
இது~$R$ ஒரு முன்கோள்
உடைய வலப்புற விதியாகவும்,
$A$ அதன் முதன்மை
!!{formula} அல்லாமலும்,
\emph{முதன்மையான} !!{formula}~$!D$
\CutCS{} இன் வலப்புற
முன்கோளின் முடிவுப்பக்கத்தில்
வருவதாகவும் உள்ள நேர்வை
மட்டுமே காட்டுகிறது.
$!D$ முன்பக்கத்தில்
வந்தாலோ
(அதாவது~$R$
இடப்புற விதியானாலோ),
அல்லது~\CutCS{} இன்
இடப்புற முன்கோளில்
வந்தாலோ வடிவம்
வேறாக இருக்கும்.
~$\Pi$,~$\Lambda$
இலுள்ள !!{formula}s
விதி~$R$ இன்
செயல்படு !!{formula}s
ஐக் குறிக்கின்றன.

இப்போது வரிசையை
மாற்றி,~\CutCS{} க்குப்
பிறகு~$R$ வரும்
!!a{proof} ஐக்
கருதுவோம்:
\[
\AxiomC{}
\RightLabel{$\pi_1'$}
\Deduce$\Gamma, \Pi \fCenter \Delta', !D, \Lambda, !A$
\AxiomC{}
\RightLabel{$\pi_3'$}
\Deduce$!A, \Gamma, \Pi \fCenter \Delta', !D, \Lambda$
\RightLabel{\CutCS}
\BinaryInf$\Gamma, \Pi \fCenter \Delta', !D, \Lambda$
\RightLabel{$R$}
\UnaryInf$\Gamma \fCenter \Delta', !D, !D$
\DisplayProof
\]
முறையே~$\pi_1$,
~$\pi_3$ இலிருந்து
~$\pi_1'$,~$\pi_3'$
பெற, !!{height}
மாறாத தளர்த்தலைப்
பயன்படுத்துகிறோம்.

\CutCS{} மற்றும்~$R$
இவற்றின் வரிசையை
மாற்றுவதால், இதனை
“வெட்டை~$R$
வழியாக மேலே
நகர்த்துதல்” என்று
அழைக்கிறோம்.
இதனால்~\CutCS{} இன்
வலப்புற முன்கோளில்
முடியும் !!{proof} இன்
உயரம் குறைகிறது;
ஆகவே வெட்டு
!!{height} உம்
குறைகிறது.
அதாவது புதிய~\CutCS{}
இன் முன்கோள்களில்
முடியும் உள்-!!{proof}s
களின் உயரங்கள்
~$\pi_1$,~$\pi_3$
இவற்றின் !!{height} களை
மீறாது எனக் கொள்ளலாம்;
வலப்புறத்தில் ஓர்
அனுமானத்தை நீக்கியதால்
வெட்டு !!{height}
குறைந்துள்ளது.
இப்போது தொகுத்தறிதல்
கருதுகோளால்~\CutCS{}
ஐ நீக்கலாம்.
இறுதியாக
$\RightR{\Contraction}$
இன் ஏற்கத்தக்கதன்மையைப்
பயன்படுத்தி
~$!D$ இன் கூடுதல்
நிகழ்வை நீக்குகிறோம்.

% SEGMENT OLP-0656-S06
\paragraph{நேர்வு C:}
~$\pi_1$,~$\pi_2$
இவற்றில் குறைந்தது
ஒன்று, ~$!A$
முதன்மையல்லாத
ஒரு முன்கோள் விதி~$R$
இல் முடிகிறது.
\CutCS{} இன் இடப்புற
அல்லது வலப்புற
முன்கோளில்~$!A$
முதன்மையல்லாதது
எதுவென்பதையும்,
ஒன்பது ஒரு முன்கோள்
விதிகளில்
(\LeftR{\lnot},
\RightR{\lnot},
\RightR{\lor},
\LeftR{\land},
\RightR{\lif},
மேலும் நான்கு
அளவையடை விதிகள்)
~$R$ எதுவென்பதையும்
பொறுத்து மொத்தம்
18 நேர்வுகள் உள்ளன.
இரண்டு எடுத்துக்காட்டுகளை
மட்டுமே பார்ப்போம்:
~$\pi_2$ இன்
கடைசி அனுமானத்தில்
~$!A$ முதன்மையல்ல;
அந்த அனுமானம்
$\RightR{\lif}$
அல்லது
\LeftR{\lexists}
இல் முடிகிறது.
மற்ற நேர்வுகள்
ஒத்தவை; அவற்றைப்
பயிற்சிகளாக விடுகிறோம்.

~$\pi$ பின்வருமாறு
முடிகிறது என்க:
\[
\AxiomC{}
\RightLabel{$\pi_1$}
\Deduce$\Gamma \fCenter \Delta', !B \lif !C, !A$
\AxiomC{}
\RightLabel{$\pi_3$}
\Deduce$!A, !B, \Gamma \fCenter \Delta', !C$
\RightLabel{$\RightR{\lif}$}
\UnaryInf$!A, \Gamma \fCenter \Delta', !B \lif !C$
\RightSubproofLabel{$\pi_2$}
\RightLabel{\CutCS}
\BinaryInf$\Gamma \fCenter \Delta', !B \lif !C$
\DisplayProof
\]
இங்கு
$\Delta = \Delta', !B \lif !C$.
~$\pi_1$ க்கு
!!{height} மாறாத தளர்த்தலை
(\olref[seq][adm]{prop:weak-G3c-adm})
பயன்படுத்தி,
$!B, \Gamma \fCenter \Delta', !B \lif !C, !C, !A$
இன் !!a{proof}~$\pi_1'$
பெறுகிறோம்
(இடப்புறத்தில்~$!B$,
வலப்புறத்தில்~$!C$
சேர்க்கப்படுகின்றன).
அதேபோல்~$\pi_3$
க்கு
\olref[seq][adm]{prop:weak-G3c-adm}
பயன்படுத்தி,
$!A, !B, \Gamma \fCenter
\Delta', !B \lif !C, !C$
இன் !!a{proof}~$\pi_3'$
பெறுகிறோம்
(வலப்புறத்தில்
$!B \lif !C$
சேர்க்கப்படுகிறது).
பின்வரும் நிறுவலுக்கு
தொகுத்தறிதல்
கருதுகோள் பொருந்தும்:
\[
\AxiomC{}
\RightLabel{$\pi_1'$}
\Deduce$!B, \Gamma \fCenter \Delta', !B \lif !C, !C, !A$
\AxiomC{}
\RightLabel{$\pi_3'$}
\Deduce$!A, !B, \Gamma \fCenter \Delta', !B \lif !C, !C$
\RightLabel{\CutCS}
\BinaryInf$!B, \Gamma \fCenter \Delta', !B \lif !C, !C$
\DisplayProof
\]
இந்த !!{proof} இன் வெட்டு
!!{formula}~$!A$;
~$\pi$ வுடன்
ஒப்பிடும்போது
வெட்டுத் தரம்
அதேயானாலும்,
வெட்டு !!{height}
குறைவு.
எனவே
\CutCS{} இன்றி
~\Log{G3c} இல்
$!B, \Gamma \fCenter \Delta', !B \lif
!C, !C$
இன் !!a{proof}~$\pi_4$
கிடைக்கிறது.
$\RightR{\lif}$
பயன்படுத்தி
$\Gamma \fCenter
\Delta', !B \lif !C, !B \lif !C$
இன் !!a{proof}
பெறுகிறோம்:
\[
\AxiomC{}
\RightLabel{$\pi_4$}
\Deduce$!B, \Gamma \fCenter \Delta', !B \lif !C, !C$
\RightLabel{$\RightR{\lif}$}
\UnaryInf$\Gamma \fCenter \Delta', !B \lif !C, !B \lif !C$
\DisplayProof
\]
அதாவது
$\Gamma \fCenter \Delta', !B \lif !C$
என்ற
$\Gamma \Sequent \Delta$
இன் !!a{proof}~$\pi'$
பெற, சுருக்கலின்
ஏற்கத்தக்கதன்மையான
\olref[seq][inv]{prop:G3c-cont-adm}
ஐப் பயன்படுத்துக.

% SEGMENT OLP-0656-S07
இப்போது~$\pi$
பின்வருமாறு
முடிகிறது என்க:
\[
\AxiomC{}
\RightLabel{$\pi_1$}
\Deduce$\lexists[x][!B(x)], \Gamma' \fCenter \Delta, !A$
\AxiomC{}
\RightLabel{$\pi_3$}
\Deduce$!A, !B(c), \Gamma' \fCenter \Delta$
\RightLabel{$\LeftR{\lexists}$}
\UnaryInf$!A, \lexists[x][!B(x)], \Gamma' \fCenter \Delta$
\RightSubproofLabel{$\pi_2$}
\RightLabel{\CutCS}
\BinaryInf$\Gamma' \fCenter \Delta$
\DisplayProof
\]
மீண்டும் பின்வரும்
நிறுவலுக்குத்
தொகுத்தறிதல்
கருதுகோளைப்
பயன்படுத்துகிறோம்:
\[
\AxiomC{}
\RightLabel{$\pi_1[!B(c) \Sequent {}]$}
\Deduce$\lexists[x][!B(x)], !B(c), \Gamma' \fCenter \Delta, !A$
\AxiomC{}
\LeftLabel{$\pi_3[\lexists[x][!B(x)] \Sequent {}]$}
\Deduce$!A, \lexists[x][!B(x)], !B(c), \Gamma' \fCenter \Delta$
\RightLabel{\CutCS}
\BinaryInf$\lexists[x][!B(x)], !B(c), \Gamma' \fCenter \Delta$
\RightLabel{\LeftR{\lexists}}
\UnaryInf$\lexists[x][!B(x)], \lexists[x][!B(x)], \Gamma' \fCenter \Delta$
\DisplayProof
\]
தனிமாறி நிபந்தனை
நிறைவேறுவதை
கவனிக்க: ~$\pi_2$
இல் உள்ள
\LeftR{\lexists}
விதியின் தனிமாறி
நிபந்தனைப்படி
$c$ ஆனது
$\lexists[x][!B(x)]$,
$\Gamma'$,
~$\Delta$ ஆகியவற்றில்
வராது. இறுதியாக
\olref[seq][inv]{prop:G3c-cont-adm}
ஐப் பயன்படுத்துக.

% SEGMENT OLP-0656-S08
\paragraph{நேர்வு D:}
~$\pi_1$,~$\pi_2$
இவற்றில் குறைந்தது
ஒன்று, ~$!A$
முதன்மையல்லாத
இரு முன்கோள்
விதி~$R$ இல்
முடிகிறது.
ஆறு நேர்வுகள்
உள்ளன.
~$\pi_2$ இன்
கடைசி அனுமானத்தில்
~$!A$ முதன்மையல்லாமல்,
அது~$\RightR{\land}$
இல் முடியும்
நேர்வைக் கருதுவோம்;
மற்ற நேர்வுகள்
ஒத்தவை.

~$\pi$ பின்வருமாறு
முடிகிறது என்க:
\[
\AxiomC{}
\RightLabel{$\pi_1$}
\Deduce$\Gamma \fCenter \Delta', !B \land !C, !A$
\AxiomC{}
\RightLabel{$\pi_3$}
\Deduce$!A, \Gamma \fCenter \Delta', !B$
\AxiomC{}
\RightLabel{$\pi_4$}
\Deduce$!A, \Gamma \fCenter \Delta', !C$
\RightLabel{$\RightR{\land}$}
\BinaryInf$!A, \Gamma \fCenter \Delta', !B \land !C$
\RightSubproofLabel{$\pi_2$}
\RightLabel{\CutCS}
\BinaryInf$\Gamma \fCenter \Delta$
\DisplayProof
\]
தொகுத்தறிதல்
கருதுகோள் பின்வரும்
!!{proof}s களுக்குப்
பொருந்தும்:
\[
\AxiomC{}
\RightLabel{$\pi_1'$}
\Deduce$\Gamma \fCenter \Delta', !B \land !C, !B, !A$
\AxiomC{}
\RightLabel{$\pi_3'$}
\Deduce$!A, \Gamma \fCenter \Delta', !B \land !C, !B$
\RightLabel{\CutCS}
\BinaryInf$\Gamma \fCenter \Delta', !B \land !C, !B$
\DisplayProof
\]
மேலும்
\[
\AxiomC{}
\RightLabel{$\pi_1''$}
\Deduce$\Gamma \fCenter \Delta', !B \land !C, !C, !A$
\AxiomC{}
\RightLabel{$\pi_4'$}
\Deduce$!A, \Gamma \fCenter \Delta', !B \land !C, !C$
\RightLabel{\CutCS}
\BinaryInf$\Gamma \fCenter \Delta', !B \land !C, !C$
\DisplayProof
\]
இங்கு~\Log{G3c}
இலுள்ள !!{proof}s
~$\pi_1'$,~$\pi_1''$
ஆகியவை~$\pi_1$
இலிருந்து !!{height}
மாறாத தளர்த்தல்
(\olref[seq][adm]{prop:weak-G3c-adm})
மூலம் பெறப்படுகின்றன:
ஒன்றில் வலப்புறம்
~$!B$ உம்,
மற்றொன்றில்~$!C$
உம் சேர்க்கப்படுகின்றன.
அதேபோல்
!!{proof}s~$\pi_3'$,
~$\pi_4'$ ஆகியவை
முறையே~$\pi_3$,
~$\pi_4$ இலிருந்து
!!{height} மாறாத
தளர்த்தல் மூலம்
பெறப்படுகின்றன;
~$\pi_3$,~$\pi_4$
இரண்டிலும்
வலப்புறம்
$!B \land !C$
சேர்க்கப்படுகிறது.

தொகுத்தறிதல்
கருதுகோள்,
$\Gamma \fCenter \Delta', !B \land !C, !B$
மற்றும்
$\Gamma \fCenter \Delta', !B \land !C, !C$
ஆகியவற்றின்
~\Log{G3c}
!!{proof}s~$\pi_5$,
~$\pi_6$ ஐத்
தருகிறது.
விதி
$\RightR{\land}$
மூலம் இவற்றை
இணைத்து,
$\Gamma \Sequent \Delta', !B \land !C, !B \land !C$
இன் !!a{proof}
பெறுகிறோம்:
\[
\AxiomC{}
\RightLabel{$\pi_5$}
\Deduce$\Gamma \fCenter \Delta', !B \land !C, !B$
\AxiomC{}
\RightLabel{$\pi_6$}
\Deduce$\Gamma \fCenter \Delta', !B \land !C, !C$
\RightLabel{\RightR{\land}}
\BinaryInf$\Gamma \fCenter \Delta', !B \land !C, !B \land !C$
\DisplayProof
\]
அதாவது
$\Gamma \fCenter \Delta', !B \land !C$
என்ற
$\Gamma \Sequent \Delta$
இன் !!a{proof}
பெற,
\olref[seq][inv]{prop:G3c-cont-adm}
வழங்கும் சுருக்கலின்
ஏற்கத்தக்கதன்மையைப்
பயன்படுத்துக.

% SEGMENT OLP-0656-S09
\paragraph{நேர்வு E: வெட்டுகளைக் குறைத்தல்.}
கடைசி நேர்வில்,
~$\pi_1$,~$\pi_2$
இரண்டும்~$!A$
முதன்மை வாய்பாடாக
உள்ள அனுமானங்களில்
முடிகின்றன.
~$!A$ இன்
சாத்தியமான ஒவ்வொரு
!!{main operator}
க்கும் ஒன்று என
ஆறு நேர்வுகள்
உள்ளன.

ஒவ்வொரு நேர்விலும்,
வெட்டு !!{formula}~$!A$
மீதான ஒரு \CutCS{}
ஐ, ~$!A$ இன்
உள்-!!{formula}s
மீதான ஒன்று அல்லது
பல வெட்டுகளாக
மாற்றுகிறோம்.
அவை~$!A$ முதன்மை
!!{formula} ஆக
உள்ள அனுமானங்களின்
செயல்படு !!{formula}s
ஆகும். எனவே இந்த
நேர்வுகள் வெட்டு
அனுமானத்தின்
“குறைப்பு” அல்லது
“மாற்றம்” எனவும்
அழைக்கப்படுகின்றன.

\begin{enumerate}
\item $!A \ident \lnot !B$
எனில்~$\pi$
பின்வருமாறு
முடிகிறது:
\[
\AxiomC{}
\RightLabel{$\pi_1'$}
\Deduce$!B, \Gamma \fCenter \Delta$
\RightLabel{\RightR{\lnot}}
\UnaryInf$\Gamma \fCenter \Delta, \lnot !B$
\LeftSubproofLabel{$\pi_1$}
\AxiomC{}
\RightLabel{$\pi_2'$}
\Deduce$\Gamma \fCenter \Delta, !B$
\RightLabel{\LeftR{\lnot}}
\UnaryInf$\lnot !B, \Gamma \fCenter \Delta$
\RightSubproofLabel{$\pi_2$}
\RightLabel{\CutCS}
\BinaryInf$\Gamma \fCenter \Delta$
\DisplayProof
\]
தொகுத்தறிதல்
கருதுகோளால்
பின்வருவதிலிருந்து
\CutCS{} ஐ
நீக்கலாம்:
\[
\AxiomC{}
\RightLabel{$\pi_2'$}
\Deduce$\Gamma \fCenter \Delta, !B$
\AxiomC{}
\RightLabel{$\pi_1'$}
\Deduce$!B, \Gamma \fCenter \Delta$
\RightLabel{\CutCS}
\BinaryInf$\Gamma \fCenter \Delta$
\DisplayProof
\]
% SEGMENT OLP-0656-S10
\item
$!A \ident (!B \lor !C)$
எனில்~$\pi$
பின்வருமாறு
முடிகிறது:
\[
\AxiomC{}
\RightLabel{$\pi_1'$}
\Deduce$\Gamma \fCenter \Delta, !B, !C$
\RightLabel{\RightR{\lor}}
\UnaryInf$\Gamma \fCenter \Delta, !B \lor !C$
\LeftSubproofLabel{$\pi_1$}
\AxiomC{}
\RightLabel{$\pi_2'$}
\Deduce$!B, \Gamma \fCenter \Delta$
\AxiomC{}
\RightLabel{$\pi_2''$}
\Deduce$!C, \Gamma \fCenter \Delta$
\RightLabel{\LeftR{\lor}}
\BinaryInf$!B \lor !C, \Gamma \fCenter \Delta$
\RightSubproofLabel{$\pi_2$}
\RightLabel{\CutCS}
\BinaryInf$\Gamma \fCenter \Delta$
\DisplayProof
\]
\CutCS{} இல் முடியும்
பின்வரும் !!{proof}
ஐக் கருதுக:
\[
\AxiomC{}
\RightLabel{$\pi_1'$}
\Deduce$\Gamma \fCenter \Delta, !B, !C$
\AxiomC{}
\RightLabel{$\pi_2'''$}
\Deduce$!C, \Gamma \fCenter \Delta, !B$
\RightLabel{\CutCS}
\BinaryInf$\Gamma \fCenter \Delta, !B$
\DisplayProof
\]
இங்கு~$\pi_2'''$
ஆனது~$\pi_2''$
இலிருந்து !!{height}
மாறாத தளர்த்தலால்
பெறப்படுகிறது.
$\depth{!C} < \depth{!B \lor !C}$
என்பதால் இதன்
வெட்டுத் தரம்
$< \cutr{\pi}$;
ஆகவே தொகுத்தறிதல்
கருதுகோள்
$\Gamma \fCenter \Delta, !B$
இன்~\Log{G3c}
இலுள்ள
!!a{proof}~$\pi_1''$
ஐத் தருகிறது.
\CutCS{} இல்
முடியும் பின்வரும்
!!{proof} ஐக்
கருதுக:
\[
\AxiomC{}
\RightLabel{$\pi_1''$}
\Deduce$\Gamma \fCenter \Delta, !B$
\AxiomC{}
\RightLabel{$\pi_2'$}
\Deduce$!B, \Gamma \fCenter \Delta$
\RightLabel{\CutCS}
\BinaryInf$\Gamma \fCenter \Delta$
\DisplayProof
\]
இதன் வெட்டுத்
தரம்
$= \depth{!B} < \depth{!B \lor !C}$;
எனவே தொகுத்தறிதல்
கருதுகோளை மீண்டும்
பயன்படுத்தி
$\Gamma \fCenter \Delta$
இன் !!a{proof}~$\pi'$
பெறுகிறோம்.
% SEGMENT OLP-0656-S11
\item $!A \ident (!B \land !C)$.
பயிற்சி.

\item $!A \ident (!B \lif !C)$
எனில்~$\pi$
பின்வருமாறு
முடிகிறது:
\[
\AxiomC{}
\RightLabel{$\pi_1'$}
\Deduce$!B, \Gamma \fCenter \Delta, !C$
\RightLabel{\RightR{\lif}}
\UnaryInf$\Gamma \fCenter \Delta, !B \lif !C$
\LeftSubproofLabel{$\pi_1$}
\AxiomC{}
\RightLabel{$\pi_2'$}
\Deduce$\Gamma \fCenter \Delta, !B$
\AxiomC{}
\RightLabel{$\pi_2''$}
\Deduce$!C, \Gamma \fCenter \Delta$
\RightLabel{\LeftR{\lif}}
\BinaryInf$!B \lif !C, \Gamma \fCenter \Delta$
\RightSubproofLabel{$\pi_2$}
\RightLabel{\CutCS}
\BinaryInf$\Gamma \fCenter \Delta$
\DisplayProof
\]
\CutCS{} இல்
முடியும் பின்வரும்
!!{proof} ஐக்
கருதுக:
\[
\AxiomC{}
\RightLabel{$\pi_1'$}
\Deduce$!B, \Gamma \fCenter \Delta, !C$
\AxiomC{}
\RightLabel{$\pi_2''$}
\Deduce$!C, \Gamma \fCenter \Delta$
\RightLabel{\CutCS}
\BinaryInf$!B, \Gamma \fCenter \Delta$
\DisplayProof
\]
இதன் வெட்டுத்
தரம்~$\pi$ இன்
வெட்டுத் தரத்தைவிடக்
குறைவு.
தொகுத்தறிதல்
கருதுகோள்
$!B, \Gamma \fCenter \Delta$
இன்
!!a{proof}~$\pi_1''$
ஐத் தருகிறது.
இப்போது~\CutCS{}
இல் முடியும்
பின்வரும்
!!{proof} ஐக்
கருதுக:
\[
\AxiomC{}
\RightLabel{$\pi_2'$}
\Deduce$\Gamma \fCenter \Delta, !B$
\AxiomC{}
\RightLabel{$\pi_1''$}
\Deduce$!B, \Gamma \fCenter \Delta$
\RightLabel{\CutCS}
\BinaryInf$\Gamma \fCenter \Delta$
\DisplayProof
\]
இதன் வெட்டுத்
தரமும்~$\pi$
இன் வெட்டுத்
தரத்தைவிடக்
குறைவு; ஆகவே
தொகுத்தறிதல்
கருதுகோள்
$\Gamma \fCenter \Delta$
இன்~\Log{G3c}
!!{proof}~$\pi'$
ஐத் தருகிறது.
% SEGMENT OLP-0656-S12
\item $!A \ident \lforall[x][!B(x)]$
எனில்~$\pi$
பின்வருமாறு
முடிகிறது:
\[
\AxiomC{}
\RightLabel{$\pi_1'$}
\Deduce$\Gamma \fCenter \Delta, !B(c)$
\RightLabel{\RightR{\lforall}}
\UnaryInf$\Gamma \fCenter \Delta, \lforall[x][!B(x)]$
\LeftSubproofLabel{$\pi_1$}
\AxiomC{}
\RightLabel{$\pi_2'$}
\Deduce$!B(t), \lforall[x][!B(x)], \Gamma \fCenter \Delta$
\RightLabel{\LeftR{\lforall}}
\UnaryInf$\lforall[x][!B(x)], \Gamma \fCenter \Delta$
\RightSubproofLabel{$\pi_2$}
\RightLabel{\CutCS}
\BinaryInf$\Gamma \fCenter \Delta$
\DisplayProof
\]
தொகுத்தறிதல்
கருதுகோளால்,
பின்வருவதிலிருந்து
\CutCS{} ஐ
நீக்கலாம்:
\[
\AxiomC{}
\RightLabel{$\pi_1''$}
\Deduce$!B(t), \Gamma \fCenter \Delta, \lforall[x][!B(x)]$
\AxiomC{}
\RightLabel{$\pi_2'$}
\Deduce$!B(t), \lforall[x][!B(x)], \Gamma \fCenter \Delta$
\RightLabel{\CutCS}
\BinaryInf$!B(t), \Gamma \fCenter \Delta$
\DisplayProof
\]
இங்கு~$\pi_1''$
ஆனது~$\pi_1$
இலிருந்து
(~$\pi_1'$ இலிருந்து
அல்ல) !!{height}
மாறாத தளர்த்தலால்
பெறப்படுகிறது.
(இந்த !!{proof}
இன் வெட்டுத்
தரம் மாறாது;
ஆனால் வெட்டு
!!{height}
குறைவாகும்.)
இதனால்
$!B(t), \Gamma \fCenter \Delta$
இன்
!!a{proof}~$\pi_2''$
கிடைக்கிறது.

!!{proof}~$\pi_1'$
இன் இறுதித் தொடரணி
$\Gamma \Sequent \Delta, !B(c)$;
இங்கு~$c$ என்பது
ஒரு
$\RightR{\lforall}$
அனுமானத்தின்
தனிமாறி.
எனவே~$c$
ஆனது
$!B(x)$,
$\Gamma$,
~$\Delta$
ஆகியவற்றில்
வராது.
!!{proof}~$\pi$
ஒழுங்கானது எனக்
கொள்கிறோம்;
ஆகவே~$c$
இந்த
$\RightR{\lforall}$
அனுமானத்திற்கு
மட்டுமே தனிமாறி;
அதன் மேலே மட்டுமே
வருகிறது.
அதாவது~$\pi_1'$
இல் மட்டுமே
வருகிறது;
~$\pi_1'$ இல்
உள்ள வேறோர்
அனுமானத்தின்
தனிமாறி அல்ல.
~$\pi_2$ இல்
$c$ வராததால்
~$t$ இலும்
அது வராது.
\olref[seq][qua]{lem:subst-regular}
மூலம்
!!{proof}~$\Subst{\pi_1'}{t}{c}$
ஆனது
$\Gamma \Sequent \Delta, !B(t)$
இன் !!a{proof}.
இப்போது
பின்வரும்
!!{proof} க்குத்
தொகுத்தறிதல்
கருதுகோளைப்
பயன்படுத்துகிறோம்:
\[
\AxiomC{}
\RightLabel{$\Subst{\pi_1'}{t}{c}$}
\Deduce$\Gamma \fCenter \Delta, !B(t)$
\AxiomC{}
\RightLabel{$\pi_2''$}
\Deduce$!B(t), \Gamma \fCenter \Delta$
\RightLabel{\CutCS}
\BinaryInf$\Gamma \fCenter \Delta$
\DisplayProof
\]
(இங்கு வெட்டு
!!{formula}~$!B(t)$;
எனவே~$\pi$
ஐவிடக் குறைந்த
வெட்டுத் தரமுள்ள
!!{proof} என்பதால்
தொகுத்தறிதல்
கருதுகோள் பொருந்தும்.)
% SEGMENT OLP-0656-S13
\item $!A \ident \lexists[x][!B(x)]$
என்ற நேர்வைப்
பயிற்சியாக விடுகிறோம்.
\end{enumerate}
\end{proof}

\begin{prob}
\olref{lem:cut-adm-G3c}
இன் நிறுவலில்
D நேர்வின்
ஒரு துணைநேர்வை
உறுதிசெய்க:
~$\pi_1$ இன்
கடைசி அனுமானத்தில்
~$!A$ முதன்மையல்ல;
அந்த அனுமானம்
$\LeftR{\lif}$
இல் முடிகிறது.
(குறிப்பு:
எதை நிறுவ வேண்டும்,
அதாவது இந்நேர்வில்
~$\pi$ எந்த
வடிவில் இருக்கும்
என்பதைத் தெளிவாகக்
கூறுவதும் இந்தப்
பயிற்சியின் ஒரு
பகுதியாகும்.)
\end{prob}

\begin{prob}
\olref{lem:cut-adm-G3c}
இன் நிறுவலில்
D நேர்வின்
ஒரு துணைநேர்வை
உறுதிசெய்க:
~$\pi_1$ இன்
கடைசி அனுமானத்தில்
~$!A$ முதன்மையல்ல;
அந்த அனுமானம்
$\LeftR{\lforall}$
இல் முடிகிறது.
\end{prob}

\begin{prob}
\olref{lem:cut-adm-G3c}
இன் நிறுவலில்
E நேர்வின்
ஒரு துணைநேர்வை
உறுதிசெய்க:
~$\pi_1$,~$\pi_2$
இரண்டின் கடைசி
அனுமானங்களிலும்
~$!A$ முதன்மையானது;
மேலும்
$!A \ident (!B \land !C)$.
\end{prob}

\begin{prob}
\olref{lem:cut-adm-G3c}
இன் நிறுவலில்
E நேர்வின்
ஒரு துணைநேர்வை
உறுதிசெய்க:
~$\pi_1$,~$\pi_2$
இரண்டின் கடைசி
அனுமானங்களிலும்
~$!A$ முதன்மையானது;
மேலும்
$!A \ident \lexists[x][!B(x)]$.
\end{prob}
% SEGMENT OLP-0656-S14
E நேர்வில்
தொகுத்தறிதல்
கருதுகோளைப்
பயன்படுத்தும்,
\CutCS{} இல்
முடியும்
!!{proof} இன்
வெட்டு !!{height},
மூல நிறுவலின்
வெட்டு !!{height} ஐவிட
அதிகமாகவும்
இருக்கலாம்
என்பதைக் கவனிக்க.
எடுத்துக்காட்டாக,
$!A \ident (!B \lif !C)$
எனில்~$\pi$
ஐ இரண்டு~\CutCS{}
அனுமானங்கள் உள்ள
!!a{proof} ஆக
மாற்றினோம்:
\[
\AxiomC{}
\RightLabel{$\pi_2'$}
\Deduce$\Gamma \fCenter \Delta, !B$
\AxiomC{}
\RightLabel{$\pi_1'$}
\Deduce$!B, \Gamma \fCenter \Delta, !C$
\AxiomC{}
\RightLabel{$\pi_2''$}
\Deduce$!C, \Gamma \fCenter \Delta$
\RightLabel{\CutCS}
\BinaryInf$!B, \Gamma \fCenter \Delta$
\RightSubproofLabel{$\pi_1''$}
\RightLabel{\CutCS}
\BinaryInf$\Gamma \fCenter \Delta$
\DisplayProof
\]
~$\pi_1'$
மற்றும்~$\pi_2''$
இடையிலுள்ள
மேலிருக்கும்~\CutCS{}
இன் வெட்டுத்
தரமும் வெட்டு
!!{height} உம்
~$\pi$ இன்
அவற்றைவிடக்
குறைவு.
ஆனால்
தொகுத்தறிதல்
கருதுகோளைப்
பயன்படுத்திய
விளைவான
!!{proof}~$\pi_1''$
க்கு இனி~$\pi$
ஐவிடக் குறைந்த
வெட்டு !!{height}
இருக்க வேண்டியதில்லை.
கீழிருக்கும்~\CutCS{}
இன் வெட்டு
!!{formula}~$!B$
என்பதால்,
அதன் வெட்டு
\emph{தரம்} மட்டும்
~$\pi$ இன்
வெட்டுத் தரத்தைவிடக்
குறைவாகவே இருக்கும்.

\olref{lem:cut-adm-G3c}
மூலம், ஓர்
!!a{proof} இல்
உச்சியிலுள்ள
ஒரே ஒரு
\CutCS{} அனுமானத்தை
நீக்கலாம்.
~\CutCS{} இல்
முடியும் உச்சி
உள்-!!{proof}s
களுக்கு அதை
ஒவ்வொன்றாகப்
பயன்படுத்தி,
ஓர் !!a{proof}
இலுள்ள எத்தனை
\CutCS{} அனுமானங்களையும்
நீக்கலாம்
எனக் காட்டலாம்.

\begin{thm}\ollabel{thm:cut-elim-G3c-topmost}
$\Log{G3c} + \CutCS$
இல் உள்ள
எந்த !!{proof}~$\pi$
ஐயும்~\Log{G3c}
இல் ஓர்
!!a{proof} ஆக
மாற்றலாம்.
\end{thm}

% SEGMENT OLP-0656-S15
\begin{proof}
  ~$\pi$ இல்
  உள்ள~\CutCS{}
  அனுமானங்களின்
  எண்ணிக்கை~$n$
  மீது தொகுத்தறிக.
  $n=0$
  எனில் செய்ய
  எதுவும் இல்லை:
  ~$\pi$ ஏற்கெனவே
  ~\Log{G3c}
  இல்
  !!a{proof}.
  இல்லையேல்,
  உச்சியிலுள்ள
  ஒரு~\CutCS{}
  அனுமானத்தில்
  முடியும்
  உள்-!!{proof}~$\pi'$
  ஐக் கருதுக.
  அதில் கடைசி
  அனுமானமாக
  ஒரே ஒரு
  \CutCS{} விதி
  உள்ளது.
  \olref{lem:cut-adm-G3c}
  ஐப் பயன்படுத்தி
  அதனை
  \CutCS{} இல்லாத
  நிறுவல்~$\pi''$
  ஆக மாற்றி,
  அதில்
  உள்-!!{proof}~$\pi'$
  க்குப் பதிலாக
  ~$\pi''$ ஐ
  வைக்கிறோம்.
  இதனால்
  $n-1$
  \CutCS{} அனுமானங்கள்
  உள்ள
  !!a{proof}
  கிடைக்கிறது;
  அதற்குத்
  தொகுத்தறிதல்
  கருதுகோள்
  பொருந்தும்.
\end{proof}

\end{document}
